% Generated from locale/tr/content/many-valued-logic/infinite-valued-logics/introduction.tex; do not hand-edit.


\olfileid[tr]{mvl}{inf}{int}

\olsection{Giriş}

Bir matrisin doğruluk değerlerinin sayısının sonlu olması gerekmez. Sonsuz
sayıda doğruluk değeri kümesi için doğal bir seçim, $0$ ile $1$ arasındaki
rasyonel sayılar kümesidir: $V_\infty = [0,1] \cap \Rat$, yani
\begin{align*}
    V_\infty & = \Setabs{\frac{n}{m}}{n,m \in \Nat \text{ ve } 0<m
    \text{ ve } n\le m}.
\intertext{Bu sonsuz doğruluk değerleri kümesini incelerken aşağıdaki alt
kümeleri de ele almak çoğu kez yararlıdır:}
V_m & = \Setabs{\frac{n}{m-1}}{n \in \Nat \text{ ve } n<m}
\intertext{Örneğin $V_5$, eşit aralıklarla dağılmış $5$ doğruluk değerinden
oluşan kümedir:}
V_5 & = \{0, \frac{1}{4}, \frac{1}{2}, \frac{3}{4}, 1\}.
\end{align*}
Bu doğruluk değeri kümelerine dayalı mantıklarda genellikle yalnızca $1$
belirlenmiştir; yani $V^+ = \{1\}$. Başka bir deyişle $1$'e (mutlak)
doğruluk, $0$'a mutlak yanlışlık rolünü veririz; fakat !!{formula}s
$V$ içindeki herhangi bir ara değeri alabilir.

Bununla birlikte $0$ ile $1$ arasındaki bütün \emph{gerçel} sayılardan
oluşan $V_{[0,1]} = [0,1]$ kümesi ya da $[0,1]$'in başka sonsuz alt
kümeleri de ele alınabilir. Bu doğruluk değerleri kümesini kullanan
mantıklara çoğu kez \emph{bulanık} denir.


